\documentclass[10pt,a4paper]{amsart}
\usepackage{amsmath,amssymb,amsthm,graphicx}
\usepackage{verbatim}
\newenvironment{reference}{\comment}{\endcomment}
%\newenvironment{reference}{}{}
\newenvironment{slogan}{\comment}{\endcomment}
\newenvironment{history}{\comment}{\endcomment}

% For commutative diagrams we use Xy-pic
\usepackage[all]{xy}

% We use 2cell for 2-commutative diagrams.
\xyoption{2cell}
% Theorem environments.
%
\theoremstyle{plain}
\newtheorem{theorem}[subsection]{Theorem}
\newtheorem{proposition}[subsection]{Proposition}
\newtheorem{lemma}[subsection]{Lemma}

\theoremstyle{definition}
\newtheorem{definition}[subsection]{Definition}
\newtheorem{example}[subsection]{Example}
\newtheorem{exercise}[subsection]{Exercise}
\newtheorem{situation}[subsection]{Situation}

\theoremstyle{remark}
\newtheorem{remark}[subsection]{Remark}
\newtheorem{remarks}[subsection]{Remarks}

\numberwithin{equation}{subsection}

% Macros
%
\def\lim{\mathop{\mathrm{lim}}\nolimits}
\def\colim{\mathop{\mathrm{colim}}\nolimits}
\def\Spec{\mathop{\mathrm{Spec}}}
\def\Hom{\mathop{\mathrm{Hom}}\nolimits}
\def\Ext{\mathop{\mathrm{Ext}}\nolimits}
\def\SheafHom{\mathop{\mathcal{H}\!\mathit{om}}\nolimits}
\def\SheafExt{\mathop{\mathcal{E}\!\mathit{xt}}\nolimits}
\def\Sch{\mathit{Sch}}
\def\Mor{\mathop{\mathrm{Mor}}\nolimits}
\def\Ob{\mathop{\mathrm{Ob}}\nolimits}
\def\Sh{\mathop{\mathit{Sh}}\nolimits}
\def\NL{\mathop{N\!L}\nolimits}
\def\CH{\mathop{\mathrm{CH}}\nolimits}
\def\proetale{{pro\text{-}\acute{e}tale}}
\def\etale{{\acute{e}tale}}
\def\QCoh{\mathit{QCoh}}
\def\Ker{\mathop{\mathrm{Ker}}}
\def\Im{\mathop{\mathrm{Im}}}
\def\Coker{\mathop{\mathrm{Coker}}}
\def\Coim{\mathop{\mathrm{Coim}}}

% Boxtimes
%
\DeclareMathSymbol{\boxtimes}{\mathbin}{AMSa}{"02}

%
% Macros for moduli stacks/spaces
%
\def\QCohstack{\mathcal{QC}\!\mathit{oh}}
\def\Cohstack{\mathcal{C}\!\mathit{oh}}
\def\Spacesstack{\mathcal{S}\!\mathit{paces}}
\def\Quotfunctor{\mathrm{Quot}}
\def\Hilbfunctor{\mathrm{Hilb}}
\def\Curvesstack{\mathcal{C}\!\mathit{urves}}
\def\Polarizedstack{\mathcal{P}\!\mathit{olarized}}
\def\Complexesstack{\mathcal{C}\!\mathit{omplexes}}
% \Pic is the operator that assigns to X its picard group, usage \Pic(X)
% \Picardstack_{X/B} denotes the Picard stack of X over B
% \Picardfunctor_{X/B} denotes the Picard functor of X over B
\def\Pic{\mathop{\mathrm{Pic}}\nolimits}
\def\Picardstack{\mathcal{P}\!\mathit{ic}}
\def\Picardfunctor{\mathrm{Pic}}
\def\Deformationcategory{\mathcal{D}\!\mathit{ef}}
\usepackage{cleveref}
\crefname{originallemma}{original lemma}{original lemmas}
\Crefname{originallemma}{Original lemma}{Original lemmas}
\makeatletter
\newlabel{decent-spaces-lemma-conditions-on-point-in-fibre-and-qf}{{0ACB}{source}{}{sourcepointer-1}{}}
\newlabel{decent-spaces-lemma-conditions-on-point-in-fibre-and-qf@cref}{{[originallemma][0][]0ACB}{[1][0][]source}{}{}{}}
\makeatother
\crefname{originallemma}{original lemma}{original lemmas}
\Crefname{originallemma}{Original lemma}{Original lemmas}
\makeatletter
\newlabel{decent-spaces-lemma-finite-over-dense-open}{{0BBC}{source}{}{sourcepointer-2}{}}
\newlabel{decent-spaces-lemma-finite-over-dense-open@cref}{{[originallemma][0][]0BBC}{[1][0][]source}{}{}{}}
\makeatother
\crefname{originalproposition}{original proposition}{original propositions}
\Crefname{originalproposition}{Original proposition}{Original propositions}
\makeatletter
\newlabel{decent-spaces-proposition-reasonable-sober}{{03K6}{source}{}{sourcepointer-3}{}}
\newlabel{decent-spaces-proposition-reasonable-sober@cref}{{[originalproposition][0][]03K6}{[1][0][]source}{}{}{}}
\makeatother
\crefname{originaltheorem}{original theorem}{original theorems}
\Crefname{originaltheorem}{Original theorem}{Original theorems}
\makeatletter
\newlabel{decent-spaces-theorem-decent-open-dense-scheme}{{086U}{source}{}{sourcepointer-4}{}}
\newlabel{decent-spaces-theorem-decent-open-dense-scheme@cref}{{[originaltheorem][0][]086U}{[1][0][]source}{}{}{}}
\makeatother
\crefname{originallemma}{original lemma}{original lemmas}
\Crefname{originallemma}{Original lemma}{Original lemmas}
\makeatletter
\newlabel{morphisms-lemma-finite-degree}{{02NX}{source}{}{sourcepointer-5}{}}
\newlabel{morphisms-lemma-finite-degree@cref}{{[originallemma][0][]02NX}{[1][0][]source}{}{}{}}
\makeatother
\crefname{originaldefinition}{original definition}{original definitions}
\Crefname{originaldefinition}{Original definition}{Original definitions}
\makeatletter
\newlabel{spaces-morphisms-definition-dominant}{{0ABL}{source}{}{sourcepointer-6}{}}
\newlabel{spaces-morphisms-definition-dominant@cref}{{[originaldefinition][0][]0ABL}{[1][0][]source}{}{}{}}
\makeatother
\crefname{originallemma}{original lemma}{original lemmas}
\Crefname{originallemma}{Original lemma}{Original lemmas}
\makeatletter
\newlabel{spaces-morphisms-lemma-compare-tr-deg}{{06RX}{source}{}{sourcepointer-7}{}}
\newlabel{spaces-morphisms-lemma-compare-tr-deg@cref}{{[originallemma][0][]06RX}{[1][0][]source}{}{}{}}
\makeatother
\title{Characterizations of a dominant morphism of integral algebraic spaces having transcendence degree zero}
\author{The Stacks Project Authors}
\date{stacks-project:spaces-over-fields.tex:local-snapshot-20261012; GNU FDL 1.2 or later}
\begin{document}
\maketitle
\noindent\textit{Editorial note.} All delivered mathematical text is the author's own, copied from the pinned Stacks Project chapter file spaces-over-fields.tex; no statement or proof step was rewritten, reordered or removed.\par\medskip
\noindent\textit{Editorial note.} The counted proof is the author's complete argument: the reduction to the schematic loci containing the generic points, the comparison of transcendence degrees and the use of finiteness implying quasi-finiteness, the identification of the points over the generic point through specializations in the fibre, and the final deduction of finiteness over a nonempty open.\par\medskip
\noindent\textit{Editorial note.} The label lemma-finite-degree names the counted statement, and the proof environment immediately following it in the source is delivered verbatim as its proof.\par\medskip
\noindent\textit{Editorial note.} The counted statement is carried by the author's own lemma environment in the chapter source.\par\medskip
\noindent\textit{Editorial note.} Pre-open gate: python operations/agent-authoring/tools/precard\_gate.py --source spaces-over-fields.tex --label lemma-finite-degree --start 603 --end 664 --local-id 28058 returned RESULT=OK (L28058 free, occupancy FREE, public interval clear). Candidate taken from sources/stacks-spaces-over-fields/candidates-halfA.json.\par\medskip
\noindent\textit{Editorial note.} The statement and proof are copied from the pinned chapter file; the cited inputs are kept as external pointers with their canonical Stacks tags.\par\medskip
\noindent\textit{Allowed standard prerequisites (editorial).}\begin{itemize}
\item Existence of a dense open schematic locus. Decent Spaces, the theorem producing a dense open subscheme and the proposition identifying reasonable spaces as decent and sober. Supplies the schematic loci Y' and X' containing the generic points, over which the morphism becomes a morphism of integral schemes.
\item The corresponding list of equivalences for morphisms of schemes. Morphisms, the lemma listing the equivalent conditions for a morphism of integral schemes: transcendence degree zero, finite residue field extension, finiteness over suitable affine opens, quasi-finiteness at the generic point and uniqueness of the point over the image. Supplies the equivalences (1),(2),(3) and the final equivalence with condition (6) in the scheme case.
\item Comparison of transcendence degrees and quasi-finiteness at a point. Morphisms of Spaces, the lemma comparing the transcendence degree of a composite and the lemma relating quasi-finiteness at a point to the absence of specializations in the fibre. Supplies the implication from (4) to (1) and the identification of the points mapping to the generic point of the target.
\item Finiteness over a dense open. Decent Spaces, the lemma producing a finite morphism over a suitable nonempty open of the target under the conditions at hand. Supplies condition (6) from the equivalent conditions (1)-(5).
\end{itemize}\medskip
\ifdefined\hypertarget\hypertarget{sourcepointer-1}{}\fi
\noindent\textit{Source reference.} Original lemma 0ACB: sources/stacks-chow/tags.txt, source lines 13049--13049. Prerequisite pointer invoked at proof line(s) 651; it names a standard fact of a sibling chapter that the counted argument uses, and every step of the counted proof built on it is retained verbatim in the delivered body.\par\medskip
\ifdefined\hypertarget\hypertarget{sourcepointer-2}{}\fi
\noindent\textit{Source reference.} Original lemma 0BBC: sources/stacks-chow/tags.txt, source lines 14240--14240. Prerequisite pointer invoked at proof line(s) 663; it names a standard fact of a sibling chapter that the counted argument uses, and every step of the counted proof built on it is retained verbatim in the delivered body.\par\medskip
\ifdefined\hypertarget\hypertarget{sourcepointer-3}{}\fi
\noindent\textit{Source reference.} Original proposition 03K6: sources/stacks-chow/tags.txt, source lines 4666--4666. Prerequisite pointer invoked at proof line(s) 631; it names a standard fact of a sibling chapter that the counted argument uses, and every step of the counted proof built on it is retained verbatim in the delivered body.\par\medskip
\ifdefined\hypertarget\hypertarget{sourcepointer-4}{}\fi
\noindent\textit{Source reference.} Original theorem 086U: sources/stacks-chow/tags.txt, source lines 10404--10404. Prerequisite pointer invoked at proof line(s) 632; it names a standard fact of a sibling chapter that the counted argument uses, and every step of the counted proof built on it is retained verbatim in the delivered body.\par\medskip
\ifdefined\hypertarget\hypertarget{sourcepointer-5}{}\fi
\noindent\textit{Source reference.} Original lemma 02NX: sources/stacks-chow/tags.txt, source lines 3496--3496. Prerequisite pointer invoked at proof line(s) 636; it names a standard fact of a sibling chapter that the counted argument uses, and every step of the counted proof built on it is retained verbatim in the delivered body.\par\medskip
\ifdefined\hypertarget\hypertarget{sourcepointer-6}{}\fi
\noindent\textit{Source reference.} Original definition 0ABL: sources/stacks-chow/tags.txt, source lines 13024--13024. Prerequisite pointer invoked at proof line(s) 608; it names a standard fact of a sibling chapter that the counted argument uses, and every step of the counted proof built on it is retained verbatim in the delivered body.\par\medskip
\ifdefined\hypertarget\hypertarget{sourcepointer-7}{}\fi
\noindent\textit{Source reference.} Original lemma 06RX: sources/stacks-chow/tags.txt, source lines 8656--8656. Prerequisite pointer invoked at proof line(s) 640; it names a standard fact of a sibling chapter that the counted argument uses, and every step of the counted proof built on it is retained verbatim in the delivered body.\par\medskip
\begin{lemma}
\label{lemma-finite-degree}
Let $S$ be a scheme. Let $X$, $Y$ be integral algebraic spaces over $S$
Let $x \in |X|$ and $y \in |Y|$ be the generic points. Let $f : X \to Y$
be locally of finite type. Assume $f$ is dominant
(Morphisms of Spaces, Definition \ref{spaces-morphisms-definition-dominant}).
The following are equivalent:
\begin{enumerate}
\item the transcendence degree of $x/y$ is $0$,
\item the extension $\kappa(x)/\kappa(y)$ (see proof) is finite,
\item there exist nonempty affine opens $U \subset X$ and $V \subset Y$
such that $f(U) \subset V$ and $f|_U : U \to V$ is finite,
\item $f$ is quasi-finite at $x$, and
\item $x$ is the only point of $|X|$ mapping to $y$.
\end{enumerate}
If $f$ is separated or if $f$ is quasi-compact, then these are
also equivalent to
\begin{enumerate}
\item[(6)] there exists a nonempty affine open $V \subset Y$ such
that $f^{-1}(V) \to V$ is finite.
\end{enumerate}
\end{lemma}

\begin{proof}
By elementary topology, we see that $f(x) = y$ as $f$ is dominant.
Let $Y' \subset Y$ be the schematic locus of $Y$ and let
$X' \subset f^{-1}(Y')$ be the schematic locus of $f^{-1}(Y')$.
By the discussion above, using
Decent Spaces, Proposition \ref{decent-spaces-proposition-reasonable-sober} and
Theorem \ref{decent-spaces-theorem-decent-open-dense-scheme},
we see that $x \in |X'|$ and $y \in |Y'|$.
Then $f|_{X'} : X' \to Y'$ is a morphism of integral schemes
which is locally of finite type. Thus we see that (1), (2), (3)
are equivalent by Morphisms, Lemma \ref{morphisms-lemma-finite-degree}.

\medskip\noindent
Condition (4) implies condition (1) by
Morphisms of Spaces, Lemma \ref{spaces-morphisms-lemma-compare-tr-deg}
applied to $X \to Y \to Y$.
On the other hand, condition (3) implies condition (4) as
a finite morphism is quasi-finite and as $x \in U$ because $x$
is the generic point. Thus (1) -- (4) are equivalent.

\medskip\noindent
Assume the equivalent conditions (1) -- (4). Suppose that
$x' \mapsto y$. Then $x \leadsto x'$ is a specialization in the
fibre of $|X| \to |Y|$ over $y$. If $x' \not = x$, then $f$ is not
quasi-finite at $x$ by Decent Spaces, Lemma
\ref{decent-spaces-lemma-conditions-on-point-in-fibre-and-qf}.
Hence $x = x'$ and (5) holds. Conversely, if (5) holds, then
(5) holds for the morphism of schemes $X' \to Y'$ (see above)
and we can use
Morphisms, Lemma \ref{morphisms-lemma-finite-degree}
to see that (1) holds.

\medskip\noindent
Observe that (6) implies the equivalent conditions (1) -- (5)
without any further assumptions on $f$. To finish the proof
we have to show the equivalent conditions (1) -- (5) imply (6).
This follows from Decent Spaces, Lemma
\ref{decent-spaces-lemma-finite-over-dense-open}.
\end{proof}

\end{document}
